\section{Discussion}
\label{sec:discussion}

\subsection{The decoupling insight (directional)}
The cleanest reading of the factorial is that the two axes play different
roles on ETTh1: \emph{routing granularity carries accuracy; KV content density
is mildly beneficial to accuracy and buys efficiency}. The control logic is
internal to the design: the DRRA $(8,4)$ cell (coarse content, fine router) matches
or beats the equal-rate cells with roughly three-fifths the attention columns
(49 vs.\ 81 complete columns), while the $(4,8)$ cell (fine content, coarse
router) is the worst cell despite having the same columns as CSA $(4,4)$.
The per-seed consistency (19/25 in the same direction for each effect), the
pre-registered H1 contrast (exact Wilcoxon $p = 0.0023$), and the non-significant
interaction ($p = 0.31$) establish this as a \emph{statistically supported sign
structure on ETTh1}; the effect sizes remain modest relative to seed noise (both
effects ${\sim}0.03$--$0.05$ vs.\ within-cell std ${\sim}0.07$), the DRRA-vs.-CSA
edge itself is directional ($p = 0.096$), and multi-dataset confirmation of the
full factorial is still future work before this can be stated as a general design
principle.

\subsection{Why might compression help on this task?}
We do not have a verified explanation for the compressed $>$ full-attention
outcome, and we explicitly do not claim one. Candidate (unverified)
explanations include: compression acting as an inductive-bias regularizer;
the block structure effectively smoothing the attention field over
multi-step correlations; and the fixed 10-epoch budget interacting with
optimization landscape differences. The cross-dataset pattern
(Section~\ref{sec:rq3}) is broadly consistent with the regularization
reading---compression helps most where adjacent timesteps are most
redundant (ETT family, solar, electricity) and slightly hurts on tasks
requiring exact in-window recall (the synthetic probes)---but it remains
descriptive evidence, not a verified mechanism. The pre-registered
hypothesis program (H1 is defined in Section~\ref{sec:rq1}) is designed to
separate the candidates, and its current status is:
\textbf{H1} (coarsening the router deteriorates accuracy more than
coarsening KV compression) directionally supported in the factorial;
\textbf{H2} (the dual-rate corner matches or beats the equal-rate cells while
using fewer attention columns), whose latency leg is supported directionally by the
systems measurements and whose accuracy leg is directional (best-mean ETTh1 cell,
DRRA-vs.-CSA $p = 0.096$; 8/12 datasets);
\textbf{H3} (dense masked attention wins in small-work regimes but loses in
large-work regimes), whose crossover shape is observed on a single machine;
\textbf{H4} (the compression advantage depends on signal properties such as
smoothness and redundancy), consistent with the cross-dataset pattern; and
\textbf{H5} (a distilled surrogate can replace attention at inference
without an accuracy loss), which is out of scope for this paper---deferred
to a separate publication (Appendix~\ref{sec:appendix-surrogate}). The
mechanistic analyses that
would separate the explanations (routing recall, lag localisation) remain
future work.

\subsection{Limitations}
\label{sec:limitations}
\begin{itemize}
  \item \textbf{Dataset breadth, not universality.} The full factorial is
        single-dataset (ETTh1); the three-way comparison spans twelve
        datasets, where the compressed advantage is signal-dependent
        (Section~\ref{sec:rq3}): wins on the ETT family, solar, and
        electricity; small losses on the synthetic probes, exchange-rate
        validation, and weather test error. No cross-scale claim is made
        (CPU-scale models only), and horizons beyond $H{=}96$ with lookbacks
        beyond $L{=}256$ are not yet evaluated.
  \item \textbf{Descriptive statistics.} Five paired seeds; per-dataset
        paired tests cannot reach significance at this sample size
        (two-sided floor $p = 0.0625$) and nothing survives Holm correction
        (Section~\ref{sec:stats}); ``best on mean'' and ``directional'' are
        used advisedly. No equivalence tests are run because no equivalence
        claims are made.
  \item \textbf{External architectures as reference points only.} DLinear
        and PatchTST were run under the frozen protocol to calibrate absolute
        error levels (Section~\ref{sec:refpoints}); they are not peer methods
        (different architectures, and DLinear has no attention), so no
        competitiveness claim is made. Efficiency-oriented peers (Informer's
        ProbSparse, iTransformer, TimesNet, SSM baselines) remain pending
        (Appendix future-work item~3).
  \item \textbf{Single-machine systems measurements.} Forward latency
        (batch and lookback scaling), MACs, memory, and training time are
        measured on the development machine (Table~\ref{tab:speed},
        Figure~\ref{fig:latency-scale} and~\ref{fig:lscale}). Cross-hardware
        results (Linux x86, GPU) and depth/width ablations are conscious
        scope exclusions, deferred to future work.
  \item \textbf{Non-IID samples.} Stride-1 windows overlap heavily; samples
        are not independent, which standard per-window significance tests
        would assume. Our paired tests are at the seed level (five
        independent runs), where the overlap does not invalidate them;
        block-bootstrap or HAC-robust corrections are deferred until
        per-window statistics are used (Section~\ref{sec:stats}).
  \item \textbf{Fixed budget, single checkpoint policy.} 10 epochs,
        epoch-10 weights; no early stopping or best-validation selection.
  \item \textbf{Single-block backbone retained.} The study keeps the minimal
        one-block backbone deliberately (the backbone, not the attention module,
        is the accuracy bottleneck at this scale, Section~\ref{sec:refpoints});
        forecast cost is reported both ways (full-layer and forecast-required,
        Section~\ref{sec:efficiency}), and a multilayer backbone remains future work.
  \item \textbf{CPU-only, fp32, small models.} Dense-with-masking
        implementation choices are CPU-oriented and may not transfer to
        GPU-flash-attention regimes.
\end{itemize}

\subsection{What this paper does not claim}
Consistent with the project's pre-registered warnings: this is not a novel
sparse-attention mechanism; it is not a state-of-the-art forecasting
system; the word ``equivalence'' is not used for any accuracy comparison;
no subquadratic claim is made beyond the structural column counts above;
and no component is claimed to be original. The contribution is the
controlled study, its measurements, and its honest boundaries.
